% Template last modified by Jake Hart; please contact course staff if you have any questions regarding using this template

\documentclass{cisXXX} % You must have the cisXXX .cls file in your project or working directory (i.e. the same directory as this document) 

\HWauthor{TRON}{} % Put your name and Penn email on this line
\HWno{1} % Enter the number of the homework you are working on
\HWcourse{Redes 2026.2} % Enter the course department and number here

\usepackage{amsmath}
\usepackage{listings}
\usepackage{xcolor}
\usepackage{dirtree}

% Configuração para boa formatação de código Python
\definecolor{codegreen}{rgb}{0,0.6,0}
\definecolor{codegray}{rgb}{0.5,0.5,0.5}
\definecolor{codepurple}{rgb}{0.58,0,0.82}
\definecolor{backcolour}{rgb}{0.95,0.95,0.92}

\lstdefinestyle{pythonstyle}{
    backgroundcolor=\color{backcolour},   
    commentstyle=\color{codegreen},
    keywordstyle=\color{magenta},
    numberstyle=\tiny\color{codegray},
    stringstyle=\color{codepurple},
    basicstyle=\ttfamily\small,
    breakatwhitespace=false,         
    breaklines=true,                 
    captionpos=b,                    
    keepspaces=true,                 
    numbers=left,                    
    numbersep=5pt,                  
    showspaces=false,                
    showstringspaces=false,
    showtabs=false,                  
    tabsize=2,
    language=Python
}

\lstset{style=pythonstyle}

\begin{document}
    \maketitle

    Esta é a documentação progressiva para o projeto-aplicativo TRON, da disciplina Redes de Computadores. Sua função é apenas descrever cada parte do código de maneira detalhada, para que cada componente tenha conhecimento completo do trabalho e o código não fique poluído com comentários complexos.

    \section{Do Gameloop}
    \subsection{Estrutura de Arquivos do Projeto}
    Os arquivos \verb|__init__.py| (mostrados como init.py) são arquivos vazios que servem somente para identificar a pasta como módulo python. Atualmente o projeto está organizado da seguinte forma:
    
    \begin{center}
    \begin{minipage}{0.7\textwidth}
    \dirtree{%
    .1 /.
    .2 TRON/ \DTcomment{Raiz do jogo}.
    .3 venv/ \DTcomment{Arquivo de ambiente virtual}.
    .3 main.py \DTcomment{Base da build/Programa cliente}.
    .3 init.py.
    .3 Entities/ \DTcomment{Definição de Classes}.
    .4 init.py.
    .4 entities.py.
    .3 Render/ \DTcomment{Funções voltadas ao uso do terminal}.
    .4 init.py.
    .4 render.py.
    .3 Lobby \DTcomment{Pasta vazia para implementar servidor/lobby}.
    }
    \end{minipage}
    \end{center}

    \subsection{Arquivo main}
            O módulo \texttt{blessed} é usado para pooling de teclado (Input não-bloqueante) com  timeout de 0.1s. Foi escolhida pois é compatível com a função \texttt{print()} nativa do Python3. O módulo \texttt{time} importa funções para usar o tempo do sistema. Aqui usamos para pausar o jogo por 1s. (Configurações não otimizadas. O jogo tem um aspecto lento e com delay de resposta considerável)

            Para rodar o jogo, é exigido um ambiente virtual \texttt{.venv} com o o módulo blessed instalado via (\texttt{pip install blessed}). Atualmente, o jogo conta com:
            \begin{itemize}
                \item Impedimento de movimento no sentido contrário da direção (Não pode dar cavalo de pau) [\textbf{\textit{Dentro da condição de mudança de direção é feita a comparação com o lado oposto.}}]
                \item Finalização do jogo no caso da presença de 1 ou menos jogadores. [\textbf{\textit{Apenas percorre o vetor de jogadores contando quantos estão vivos (isAlive = True)}}]
                \item 4 cores para melhor destaque Vermelho, Verde, Amarelo e Azul
                \item Colisão dos jogadores com extremos da arena ("/"), com rastro do proprio jogo e rastro dos demais jogadores em todas as direções. [\textbf{\textit{A verficação é simples: move-se o jogador e logo, em seguida, é feita a comparação com o valor do tile. Se for um tile preenchido \texttt{jogador.isAlive = False}}}]
                \item Geração de rastro da mesma cor do jogador. Indicado por "*" colorido [\textit{\textbf{Atualiza a antiga posição do jogador para "*" e depois o jogador é movimentado.}}]
            \end{itemize}
    
    \begin{lstlisting}
    from .Entities.entities import Cor, Arena, Jogador, Direcao
    from .Render import render
    import time
    from blessed import Terminal
    
    
    def main():
    
        # Deteccao e instanciacao da lista de jogadores (Lobby/Servidor) #
    
        LARGURA = 50
        ALTURA = 50
        arena = Arena(jogadores, LARGURA, ALTURA)
    
        # GAME LOOP
        term = Terminal()
        with term.cbreak():
            while True:
                # Pooling de teclado
                if jogador1.get_isAlive():
                    tecla = term.inkey(timeout=0.1)
                    if tecla.lower() == "w" and jogador1.direcao != Direcao.DOWN:
                        jogador1.set_direcao(Direcao.UP)
                    elif tecla.lower() == "a" and jogador1.direcao != Direcao.RIGHT:
                        jogador1.set_direcao(Direcao.LEFT)
                    elif tecla.lower() == "s" and jogador1.direcao != Direcao.UP:
                        jogador1.set_direcao(Direcao.DOWN)
                    elif tecla.lower() == "d" and jogador1.direcao != Direcao.LEFT:
                        jogador1.set_direcao(Direcao.RIGHT)
    
                # Renderizacao
                render.limpar_tela()
                arena.show()
    
                vivos = 0
                for j in jogadores:
                    if j.get_isAlive():
                        vivos += 1
    
                if vivos <= 1:
                    print("Jogo encerrou")
                    break
    
                arena.atualiza()
                time.sleep(0.5)
    
    
    if __name__ == "__main__":
        main()
    \end{lstlisting}
    \subsection{Das classes}
    As cores são definidas usando código ANSI para cor. Útil para casos simples como esse
    \begin{lstlisting}
        class Cor:
            RED = "\033[31m"
            GREEN = "\033[32m"
            YELLOW = "\033[33m"
            BLUE = "\033[34m"
            RESET = "\033[0m"
    \end{lstlisting}

    A direção é uma tupla com dois valores: incremento do eixo X e incremento do eixo Y.
    \begin{lstlisting}
        class Direcao:
            UP = (0, -1)
            DOWN = (0, 1)
            LEFT = (-1, 0)
            RIGHT = (1, 0)
            NEUTRO = (0, 0)
    \end{lstlisting}

    Um jogador apresenta cor prória, coordenadas e direcao (Questiona-se colocar o endereço e porta dentro da classe).
    De resto tem um getter, um setter e função move(), que é autoexplicativa.
    \begin{lstlisting}
        class Jogador:
            def __init__(self, pos_x: int, pos_y: int, cor: Cor, direcao: Direcao):
                self.pos_x = pos_x
                self.pos_y = pos_y
                self.cor = cor
                self.isAlive = True
                self.direcao = direcao
        
            def get_isAlive(self):
                return self.isAlive
        
            def set_direcao(self, direcao):
                self.direcao = direcao
        
            def move(self):
                dx, dy = self.direcao
                self.pos_x += dx
                self.pos_y += dy
    \end{lstlisting}

    Por fim, a classe Arena descreve o espaço de jogo. Trata-se de uma matriz de strings que salva a cor e o símbolo de cada tile dentro da string da célula. Conta com duas funções show() que printa a matriz e a funciona atualiza que movimenta todos jogadores e coloca seus rastros.
    \begin{lstlisting}
    class Arena:
        def __init__(self, jogadores, largura, altura):
            self.jogadores = jogadores
            self.LARGURA = largura
            self.ALTURA = altura
            self.grid = [[" " for _ in range(self.LARGURA)] for _ in range(self.ALTURA)]
            for i in range(self.ALTURA):
                for j in range(self.LARGURA):
                    if i == 0 or j == 0 or i == self.ALTURA-1 or j == self.LARGURA-1:
                        self.grid[i][j] = "/"
    
        def show(self):
            for j in self.jogadores:
                if(j.get_isAlive()):
                    self.grid[j.pos_y][j.pos_x] = f"{j.cor}&{Cor.RESET}"
    
            for i in range(self.ALTURA):
                for j in range(self.LARGURA):
                    print(self.grid[i][j], end=" ")
                print()
    
        def atualiza(self):
            for j in self.jogadores:
                if(j.get_isAlive()):
                    self.grid[j.pos_y][j.pos_x] = f"{j.cor}*{Cor.RESET}"
                    j.move()
                    if "/" in self.grid[j.pos_y][j.pos_x] or "*" in self.grid[j.pos_y][j.pos_x]:
                        j.isAlive = False
    \end{lstlisting}
    
    \subsection{Da renderização}
    É um módulo com só uma função no momento - \verb|limpa_tela()|, sua utilidade é auto-explicativa. Foi adaptada usando os módulos \texttt{os} e \texttt{subprocess} para funcionar em qualquer sistema operacional. (Ainda não foi testado no Windows)
    \begin{lstlisting}
    import os
    import subprocess
    
    def limpar_tela():
        subprocess.run("cls" if os.name=="nt" else "clear")
    
    \end{lstlisting}

    \begin{figure}
        \centering
        \includegraphics[width=0.5\linewidth]{}
        \caption{Caption}
        \label{fig:placeholder}
    \end{figure}

    \section{Da Network}

\end{document}